% =====================================================================
%  MSACL DS301 Deep Learning · Lab 4 · Code Hint Sheet
%  Multiple-choice code options for the three fill-in cells.
%  Style matches labs/handouts/lab02_hints.tex and lab03_hints.tex.
% =====================================================================
\documentclass[11pt]{article}

\usepackage[letterpaper, margin=0.55in, top=0.55in, bottom=0.4in]{geometry}
\usepackage{amsmath, amssymb}
\usepackage{xcolor}
\usepackage{fancyhdr}
\usepackage{enumitem}

\pagestyle{fancy}
\fancyhf{}
\fancyhead[L]{\small\textbf{MSACL DS301 --- Deep Learning}}
\fancyhead[R]{\small\textbf{Lab 4: Code Hint Sheet}}
\fancyfoot[C]{\thepage}
\renewcommand{\headrulewidth}{0.4pt}

\setlength{\parindent}{0pt}
\setlength{\parskip}{1.5pt}
\newcommand{\opt}[1]{\texttt{#1}}
\setlist[itemize]{leftmargin=2.2em, itemsep=0pt, topsep=0pt, parsep=0pt}

\begin{document}
\small

\begin{center}
{\Large \textbf{Lab 4 --- Code Hint Sheet}}\\[2pt]
{\normalsize A VAE for Spectra QC, Anomaly Detection \& Generation}
\end{center}

\vspace{-2pt}
\begin{center}
\fbox{\parbox{0.92\textwidth}{\small
Stuck on a \textbf{YOUR TURN} cell? Each blank you need is below with a few options.
Pick one, type it in, and \textbf{run the cell} --- the built-in \texttt{assert} checks will tell you
whether it is right. Only the three marked cells need editing; run everything else as-is.
The encoder, decoder, and the reparameterization trick are already written for you.
}}
\end{center}

\vspace{2pt}\hrule\vspace{4pt}

\textbf{Step 3 --- the bottleneck (latent) size.}\; {\small The narrow layer every spectrum must
squeeze through (Lecture 11). Keep it \emph{small} so unfamiliar spectra can't be rebuilt --- that is
what makes their reconstruction error spike. Value \opt{2} also lets us plot the latent space directly.}

\textbf{Blank 1 --- \texttt{LATENT\_DIM}.}
\begin{itemize}
  \item[(a)] \opt{LATENT\_DIM = 2}
  \item[(b)] \opt{LATENT\_DIM = 256}
  \item[(c)] \opt{LATENT\_DIM = 1}
\end{itemize}
{\small (a) is right --- small enough to force real compression, and exactly \opt{2} so the 2-D latent
plot in Step 9 works (the check requires \opt{LATENT\_DIM == 2}). (b) is \textbf{too big}: a wide
bottleneck lets the network copy \emph{any} input straight through, so even weird spectra reconstruct
well and the \textbf{anomaly detector stops working}. (c) is \textbf{too tight} and can't be plotted in
2-D --- both (b) and (c) fail the \opt{== 2} check.}

\vspace{3pt}\hrule\vspace{4pt}

\textbf{Step 5 --- the combined VAE loss.}\; {\small A VAE is trained on \emph{two terms added
together} (Lecture 11's recipe slide): the \opt{recon\_loss} (rebuild the spectrum accurately) plus the
\opt{kl\_loss} (keep the latent cloud tidy), with the KL weighted by \opt{beta}. Both terms are already
computed for you --- write the one line that combines them.}

\textbf{Blank 2 --- \texttt{loss}.}
\begin{itemize}
  \item[(a)] \opt{loss = recon\_loss + beta * kl\_loss}
  \item[(b)] \opt{loss = recon\_loss}
  \item[(c)] \opt{loss = recon\_loss - beta * kl\_loss}
\end{itemize}
{\small (a) is the VAE objective: reconstruction \emph{plus} the weighted KL ``tidiness'' term. (b)
drops the KL entirely --- that is just a plain autoencoder, and the check (loss must equal
\opt{recon\_loss + BETA*kl\_loss}) fails. (c) \emph{subtracts} the KL, which rewards a messy latent and
also fails the check --- and would blow training up. (Turning \opt{BETA} way \emph{up} instead makes
reconstructions blurry --- tidiness beats accuracy.)}

\vspace{3pt}\hrule\vspace{4pt}

\textbf{Step 7 --- the anomaly threshold.}\; {\small Turn reconstruction error into a flag. Measure the
error on the \emph{normal training} spectra, then flag anything well above that. A high percentile ---
the \textbf{95th} --- expects to false-alarm on only $\sim$5\% of normal spectra. This is a
sensitivity/specificity dial (Lecture 10 callback).}

\textbf{Blank 3 --- \texttt{THRESHOLD}.}
\begin{itemize}
  \item[(a)] \opt{THRESHOLD = np.percentile(train\_errors, 95)}
  \item[(b)] \opt{THRESHOLD = 0}
  \item[(c)] \opt{THRESHOLD = train\_errors.max()}
\end{itemize}
{\small (a) is the 95th-percentile cut --- it lands strictly \emph{inside} the range of training errors,
which is exactly what the check requires. (b) \opt{= 0} sits below every error, so it flags
\textbf{everything} as an anomaly (all false alarms) and fails the ``inside the range'' check at the low
end. (c) \opt{.max()} sits \emph{at} the very top, so it flags \textbf{almost nothing} (misses every
anomaly) --- the opposite failure, and it too fails the check, which requires the threshold to sit
strictly \emph{below} the max. You can raise the percentile for fewer false alarms or lower it to catch
more; \opt{99} also passes (still below the max).}

\vspace{4pt}\hrule\vspace{3pt}
{\small\itshape Answers self-check in the notebook: if the \texttt{assert} cells pass ---
\opt{LATENT\_DIM == 2}, the loss equals \opt{recon\_loss + BETA*kl\_loss} (a finite scalar, KL $\ge 0$),
and the threshold falls strictly between the min and max training error --- you chose right. Then read the
histograms \textbf{honestly}: resistant and normal spectra \emph{overlap} --- reconstruction error is a
\textbf{QC / novelty} signal, not a resistance test. The payoff is the smooth VAE latent map (vs. the
scattered plain-AE one) and the decoded spectrum morph.}

\end{document}
